\subsubsection{Resolución de problemas}


	\ejemplosboxed[4]{Resuelve los siguientes problemas:

		\begin{multicols}{2}
			\begin{parts}\footnotesize%
				\part Calcula la altura de un prisma que tiene como área de la base 6 m$^2$ y 99 m$^3$ de capacidad.

				\begin{solutionbox}{2cm}
					Ya que el volumen de un prisma es: $V=A_b\cdot h$, entonces la altura del prisma es: \[h=\dfrac{V}{A_b}=\dfrac{99}{6}=16.5\text{m}\]
				\end{solutionbox}

				% \part ¿Cuál es el perímetro de un campo de fútbol que mide 95.12 metros de largo y 45.27 metros de ancho?

				% \begin{solutionbox}{1.8cm}
				% 	El perímetro de un rectángulo es: $P=2(l+a)$ entonces el perímetro del campo de fútbol es: \[P=2(95.12+45.27)=280.78\text{m}\]
				% \end{solutionbox}

				\columnbreak%

				% \part Calcula la altura de un prisma que tiene como área de la base 8 m$^2$ y 144 m$^3$ de capacidad.

				% \begin{solutionbox}{2cm}
				% 	Ya que el volumen de un prisma es: $V=A_b\cdot h$, entonces la altura del prisma es: \[h=\dfrac{V}{A_b}=\dfrac{144}{8}=18\text{m}\].
				% \end{solutionbox}				

				\part Ricardo quiere poner una barda alrededor de un terreno pentagonal que mide 15 metros por lado. ¿Cuánta barda necesitará Ricardo para poner barda en todo el terreno?

				\begin{solutionbox}{1.8cm}
					Se sabe que el perímetro de un pentágono es: $P=5l$, entonces el perímetro del terreno es: \[P=5(15)=75\text{m}\]
				\end{solutionbox}
			\end{parts}
		\end{multicols}
	}

	\questionboxed[4]{Resuelve los siguientes problemas:

		\begin{multicols}{2}
			\begin{parts}\footnotesize%
				% \part Calcula la altura de un prisma que tiene como área de la base 6 m$^2$ y 99 m$^3$ de capacidad.

				% \begin{solutionbox}{2cm}
				% 	Ya que el volumen de un prisma es: $V=A_b\cdot h$, entonces la altura del prisma es: \[h=\dfrac{V}{A_b}=\dfrac{99}{6}=16.5\text{m}\]
				% \end{solutionbox}

				\part ¿Cuál es el perímetro de un campo de fútbol que mide 95.12 metros de largo y 45.27 metros de ancho?

				\begin{solutionbox}{1.8cm}
					El perímetro de un rectángulo es: $P=2(l+a)$ entonces el perímetro del campo de fútbol es: \[P=2(95.12+45.27)=280.78\text{m}\]
				\end{solutionbox}

				\columnbreak%

				\part Calcula la altura de un prisma que tiene como área de la base 8 m$^2$ y 144 m$^3$ de capacidad.

				\begin{solutionbox}{2cm}
					Ya que el volumen de un prisma es: $V=A_b\cdot h$, entonces la altura del prisma es: \[h=\dfrac{V}{A_b}=\dfrac{144}{8}=18\text{m}\].
				\end{solutionbox}				

				% \part Ricardo quiere poner una barda alrededor de un terreno pentagonal que mide 15 metros por lado. ¿Cuánta barda necesitará Ricardo para poner barda en todo el terreno?

				% \begin{solutionbox}{1.8cm}
				% 	Se sabe que el perímetro de un pentágono es: $P=5l$, entonces el perímetro del terreno es: \[P=5(15)=75\text{m}\]
				% \end{solutionbox}
			\end{parts}
		\end{multicols}
	}
	
	